\documentclass[dvipsnames]{article}
\usepackage{geometry}
\geometry{a3paper, landscape, margin=2in}
\usepackage{url}
\usepackage{multicol}
\usepackage{esint}
\usepackage{amsfonts}
\usepackage{tikz}
\usetikzlibrary{decorations.pathmorphing}
\usepackage{amsmath,amssymb}
\usepackage{enumitem}
\usepackage{colortbl}
\usepackage{mathtools}
\usepackage{ mathrsfs }	

\usepackage{ gensymb }

\usepackage[activate={true,nocompatibility},final,tracking=true,kerning=true,spacing=true,factor=1100,stretch=10,shrink=10]{microtype}
\makeatletter
\setlist[itemize]{noitemsep, topsep=0pt}
\usepackage{quiver}
\newcommand*\bigcdot{\mathpalette\bigcdot@{.5}}
\newcommand*\bigcdot@[2]{\mathbin{\vcenter{\hbox{\scalebox{#2}{$\m@th#1\bullet$}}}}}
\makeatother
%Sheet made by Boris, Template is by Drew Ulick https://de.overleaf.com/articles/130-cheat-sheet/ntwtkmpxmgrp
\usepackage[english]{babel}
\usepackage{lmodern}
\usepackage{bm}
	\definecolor{cobalt}{rgb}{0.0, 0.28, 0.67}
\advance\topmargin-1.2in
\advance\textheight3in
\advance\textwidth3in
\advance\oddsidemargin-1.5in
\advance\evensidemargin-1.5in
\parindent0pt
\parskip2pt
\newcommand{\hr}{\centerline{\rule{3.5in}{1pt}}}
\newcommand{\R}{\mathbb{R}}
\newcommand{\Q}{\mathbb{Q}}
\newcommand{\C}{\mathbb{C}}
\newcommand{\Z}{\mathbb{Z}}
\newcommand{\N}{\mathbb{N}}
\newcommand{\D}{\mathbb{D}}
\newcommand{\F}{\mathbb{F}}
% \colorbox[HTML]{e4e4e4}{\makebox[\textwidth-2\fboxsep][l]{texto}}
\begin{document}
\definecolor{orange}{RGB}{255, 115,0}
\begin{center}{\fontsize{35}{60}{\textcolor{DarkOrchid}{\textbf{Algebraic Topology Cheat Sheet}}}}\\
\end{center}

\tikzset{every loop/.style={min distance=5mm,in=20,out=90,looseness=1}}
\tikzstyle{mybox} = [draw=DarkOrchid, fill=white, very thick,
rectangle, rounded corners, inner sep=10pt, inner ysep=10pt]
\tikzstyle{fancytitle} =[fill=DarkOrchid, text=white, font=\bfseries]

\begin{multicols*}{3}

% ------------------------------------------------------
% Start of Sheet
% ------------------------------------------------------
%------------ Topology ---------------
\begin{tikzpicture}
	\node [mybox] (box){%
		\begin{minipage}{0.3\textwidth}
			A topology on a set $X$ is the tuple $(X,T)$ where $T$ is a collection of (open) subsets of $X$ such that 
			\begin{itemize}
				\item $\emptyset \in T, X \in T$.
				\item Closed under unions.
				\item Closed under finite intersections.
			\end{itemize}
		\end{minipage}
	};
	%------------ Topology Header ---------------------
	\node[fancytitle, right=10pt] at (box.north west) {Topology};
\end{tikzpicture}
%------------ Constructions on topologies ---------------
\begin{tikzpicture}
	\node [mybox] (box){%
		\begin{minipage}{0.3\textwidth}
			\small
			Let \(X\) be a topological space.
			
			If \(Y\subseteq X\), the \textbf{subspace topology} on \(Y\) is
			\[
			\mathcal T_Y=\{U\cap Y:U\in \mathcal T_X\}.
			\]
			Equivalently, for the inclusion \(i:Y\hookrightarrow X\), a map \(h:Z\to Y\) is continuous iff \(i\circ h:Z\to X\) is continuous.
			\[
			\begin{tikzcd}[column sep=small,row sep=small]
				Z \arrow[r,"h"] \arrow[dr,"g=i\circ h"'] & Y \arrow[d,hook,"i"] \\
				& X
			\end{tikzcd}
			\]
			
			If \(q:X\twoheadrightarrow S\) is surjective, the \textbf{quotient topology} on \(S\) is defined by
			\[
			U\subseteq S \text{ open }\iff q^{-1}(U)\subseteq X \text{ open}.
			\]
			Equivalently, a map \(h:S\to Z\) is continuous iff \(h\circ q:X\to Z\) is continuous.
			\[
			\begin{tikzcd}[column sep=small,row sep=small]
				X \arrow[r,two heads,"q"] \arrow[dr,"g=h\circ q"'] & S \arrow[d,"h"] \\
				& Z
			\end{tikzcd}
			\]
			
			For a family \(\{X_a\}\), the \textbf{product topology} on \(\prod_a X_a\) is characterized by: \(h:Z\to \prod_a X_a\) is continuous iff each \(\pi_a\circ h:Z\to X_a\) is continuous.
			\[
			\begin{tikzcd}[column sep=small,row sep=small]
				& \prod_a X_a \arrow[d,"\pi_a"] \\
				Z \arrow[ur,"h"] \arrow[r,"g_a=\pi_a\circ h"'] & X_a
			\end{tikzcd}
			\]
			
			The \textbf{coproduct topology} on \(\coprod_a X_a\) is dual: \(h:\coprod_a X_a\to Z\) is continuous iff each \(h\circ \iota_a:X_a\to Z\) is continuous.
			\[
			\begin{tikzcd}[column sep=small,row sep=small]
				X_a \arrow[r,"\iota_a"] \arrow[dr,"h\circ \iota_a"'] & \coprod_a X_a \arrow[d,"h"] \\
				& Z
			\end{tikzcd}
			\]
		\end{minipage}
	};
	\node[fancytitle, right=10pt] at (box.north west) {Constructions on topologies};
\end{tikzpicture}


%------------ Deformation retraction ---------------
\begin{tikzpicture}
\node [mybox] (box){%
    \begin{minipage}{0.3\textwidth}
	  Let \(A\subseteq X\). A \textbf{deformation retraction} of \(X\) onto \(A\) is a homotopy \(H:X\times I\to X\) such that \(H(x,0)=x\), \(H(x,1)\in A\), and \(H(a,t)=a\) for all \(a\in A\), \(t\in I\). Equivalently, \(H_0=\mathrm{id}_X\), \(H_1=r\) for a retraction \(r:X\to A\), and \(H_t|_A=\mathrm{id}_A\) for all \(t\in I\).
  	
    \end{minipage}
};
%------------ Deformation retraction Header ---------------------
\node[fancytitle, right=10pt] at (box.north west) {Deformation retraction};
\end{tikzpicture}

%------------ Homotopy ---------------
\begin{tikzpicture}
	\node [mybox] (box){%
		\begin{minipage}{0.3\textwidth}
		Let \(f_0,f_1:X\to Y\) be continuous. A \textbf{homotopy} from \(f_0\) to \(f_1\) is a continuous map \(H:X\times I\to Y\) such that \(H(x,0)=f_0(x)\) and \(H(x,1)=f_1(x)\). Write \(f_0\simeq f_1\).
		
		Spaces \(X\) and \(Y\) are \textbf{homotopy equivalent} if there exist continuous maps \(f:X\to Y\) and \(g:Y\to X\) such that \(g\circ f\simeq \mathrm{id}_X\) and \(f\circ g\simeq \mathrm{id}_Y\).
		
		Homeomorphic \(\Rightarrow\) homotopy equivalent, but not conversely.
		
		A space \(X\) is \textbf{contractible} if \(\mathrm{id}_X\) is homotopic to a constant map.
		
		A map \(f:X\to Y\) is \textbf{null-homotopic} if \(f\simeq c\) for some constant map \(c:X\to Y\).
		\end{minipage}
	};
	%------------ Homotopy Header ---------------------
	\node[fancytitle, right=10pt] at (box.north west) {Homotopy};
\end{tikzpicture}


%------------ CW Complex ---------------
\begin{tikzpicture}
	\node [mybox] (box){%
		\begin{minipage}{0.3\textwidth}
			A \textbf{CW complex} is a space \(X=\bigcup_{n\ge 0}X^n\) such that \(X^0\) is a discrete set and, for each \(n\ge 1\), \(X^n\) is obtained from \(X^{n-1}\) by attaching \(n\)-cells \(e_\alpha^n\cong D^n\) along maps \(\varphi_\alpha:S^{n-1}\to X^{n-1}\). Equivalently, \(X^n\) is defined by the pushout
			\[
			\begin{tikzcd}[column sep=small,row sep=small]
				\coprod_\alpha S^{n-1}_\alpha & X^{n-1} \\
				\coprod_\alpha D^n_\alpha & X^n
				\arrow["{\coprod \varphi_\alpha}", from=1-1, to=1-2]
				\arrow[hook, from=1-1, to=2-1]
				\arrow[from=1-2, to=2-2]
				\arrow[from=2-1, to=2-2]
				\arrow["\lrcorner"{anchor=center, pos=0.125, rotate=180}, draw=none, from=2-2, to=1-1]
			\end{tikzcd}
			\]
			and \(X^n\setminus X^{n-1}=\coprod_\alpha e_\alpha^n\). The spaces \(X^n\) are the \(n\)-skeleta of \(X\).
		\end{minipage}
	};
	\node[fancytitle, right=10pt] at (box.north west) {CW Complex};
\end{tikzpicture}
%------------ Constructions on CW complexes ---------------
\begin{tikzpicture}
	\node [mybox] (box){%
		\begin{minipage}{0.3\textwidth}
			
		\end{minipage}
	};
	%------------ Constructions on CW complexes Header ---------------------
	\node[fancytitle, right=10pt] at (box.north west) {Constructions on CW complexes};
\end{tikzpicture}

%------------ Examples of CW structures ---------------
\begin{tikzpicture}
	\node [mybox] (box){%
		\begin{minipage}{0.3\textwidth}
			
		\end{minipage}
	};
	%------------ Examples of CW structures Header ---------------------
	\node[fancytitle, right=10pt] at (box.north west) {Examples of CW Structures};
\end{tikzpicture}

%------------ Paths and loops ---------------
\begin{tikzpicture}
	\node [mybox] (box){%
		\begin{minipage}{0.3\textwidth}
			
		\end{minipage}
	};
	%------------ Paths and loops Header ---------------------
	\node[fancytitle, right=10pt] at (box.north west) {Paths and loops};
\end{tikzpicture}

%------------ Fundamental groupoid ---------------
\begin{tikzpicture}
	\node [mybox] (box){%
		\begin{minipage}{0.3\textwidth}
			
		\end{minipage}
	};
	%------------ Fundamental groupoid Header ---------------------
	\node[fancytitle, right=10pt] at (box.north west) {Fundamental groupoid};
\end{tikzpicture}

%------------ van Kampen theorem for groupoids ---------------
\begin{tikzpicture}
	\node [mybox] (box){%
		\begin{minipage}{0.3\textwidth}
			
		\end{minipage}
	};
	%------------ van Kampen theorem for groupoids Header ---------------------
	\node[fancytitle, right=10pt] at (box.north west) {van Kampen theorem for groupoids};
\end{tikzpicture}

%------------ Fundamental group ---------------
\begin{tikzpicture}
	\node [mybox] (box){%
		\begin{minipage}{0.3\textwidth}
			
		\end{minipage}
	};
	%------------ Fundamental group Header ---------------------
	\node[fancytitle, right=10pt] at (box.north west) {Fundamental group};
\end{tikzpicture}

%------------ Covering spaces ---------------
\begin{tikzpicture}
	\node [mybox] (box){%
		\begin{minipage}{0.3\textwidth}
			
		\end{minipage}
	};
	%------------ Covering spaces Header ---------------------
	\node[fancytitle, right=10pt] at (box.north west) {Covering spaces};
\end{tikzpicture}

%------------ Lifting and deck transformations ---------------
\begin{tikzpicture}
	\node [mybox] (box){%
		\begin{minipage}{0.3\textwidth}
			
		\end{minipage}
	};
	%------------ Lifting and deck transformations Header ---------------------
	\node[fancytitle, right=10pt] at (box.north west) {Lifting and deck transformations};
\end{tikzpicture}

%------------ Classification of covering spaces ---------------
\begin{tikzpicture}
	\node [mybox] (box){%
		\begin{minipage}{0.3\textwidth}
			
		\end{minipage}
	};
	%------------ Classification of covering spaces Header ---------------------
	\node[fancytitle, right=10pt] at (box.north west) {Classification of covering spaces};
\end{tikzpicture}

%------------ Higher homotopy groups ---------------
\begin{tikzpicture}
	\node [mybox] (box){%
		\begin{minipage}{0.3\textwidth}
			
		\end{minipage}
	};
	%------------ Higher homotopy groups Header ---------------------
	\node[fancytitle, right=10pt] at (box.north west) {Higher homotopy groups};
\end{tikzpicture}

%------------ Eckmann-Hilton argument ---------------
\begin{tikzpicture}
	\node [mybox] (box){%
		\begin{minipage}{0.3\textwidth}
			
		\end{minipage}
	};
	%------------ Eckmann-Hilton argument Header ---------------------
	\node[fancytitle, right=10pt] at (box.north west) {Eckmann-Hilton argument};
\end{tikzpicture}

%------------ Whitehead product ---------------
\begin{tikzpicture}
	\node [mybox] (box){%
		\begin{minipage}{0.3\textwidth}
			
		\end{minipage}
	};
	%------------ Whitehead product Header ---------------------
	\node[fancytitle, right=10pt] at (box.north west) {Whitehead product};
\end{tikzpicture}

%------------ Hurewicz theorem ---------------
\begin{tikzpicture}
	\node [mybox] (box){%
		\begin{minipage}{0.3\textwidth}
			
		\end{minipage}
	};
	%------------ Hurewicz theorem Header ---------------------
	\node[fancytitle, right=10pt] at (box.north west) {Hurewicz theorem};
\end{tikzpicture}

%------------ Whitehead theorem ---------------
\begin{tikzpicture}
	\node [mybox] (box){%
		\begin{minipage}{0.3\textwidth}
			
		\end{minipage}
	};
	%------------ Whitehead theorem Header ---------------------
	\node[fancytitle, right=10pt] at (box.north west) {Whitehead theorem};
\end{tikzpicture}

%------------ Fibrations and cofibrations ---------------
\begin{tikzpicture}
	\node [mybox] (box){%
		\begin{minipage}{0.3\textwidth}
			
		\end{minipage}
	};
	%------------ Fibrations and cofibrations Header ---------------------
	\node[fancytitle, right=10pt] at (box.north west) {Fibrations and cofibrations};
\end{tikzpicture}

%------------ Postnikov towers ---------------
\begin{tikzpicture}
	\node [mybox] (box){%
		\begin{minipage}{0.3\textwidth}
			
		\end{minipage}
	};
	%------------ Postnikov towers Header ---------------------
	\node[fancytitle, right=10pt] at (box.north west) {Postnikov towers};
\end{tikzpicture}

%------------ Obstruction theory ---------------
\begin{tikzpicture}
	\node [mybox] (box){%
		\begin{minipage}{0.3\textwidth}
			
		\end{minipage}
	};
	%------------ Obstruction theory Header ---------------------
	\node[fancytitle, right=10pt] at (box.north west) {Obstruction theory};
\end{tikzpicture}

%------------ Spectral sequences ---------------
\begin{tikzpicture}
	\node [mybox] (box){%
		\begin{minipage}{0.3\textwidth}
			
		\end{minipage}
	};
	%------------ Spectral sequences Header ---------------------
	\node[fancytitle, right=10pt] at (box.north west) {Spectral sequences};
\end{tikzpicture}

%------------ Simplicial, singular and cellular homology ---------------
\begin{tikzpicture}
	\node [mybox] (box){%
		\begin{minipage}{0.3\textwidth}
			
		\end{minipage}
	};
	%------------ Simplicial, singular and cellular homology Header ---------------------
	\node[fancytitle, right=10pt] at (box.north west) {Simplicial, singular and cellular homology};
\end{tikzpicture}

%------------ Homology of Pairs ---------------
\begin{tikzpicture}
	\node [mybox] (box){%
		\begin{minipage}{0.3\textwidth}
			
		\end{minipage}
	};
	%------------ Homology of Pairs Header ---------------------
	\node[fancytitle, right=10pt] at (box.north west) {Homology of Pairs};
\end{tikzpicture}

%------------ Reduced homology ---------------
\begin{tikzpicture}
	\node [mybox] (box){%
		\begin{minipage}{0.3\textwidth}
			
		\end{minipage}
	};
	%------------ Reduced homology Header ---------------------
	\node[fancytitle, right=10pt] at (box.north west) {Reduced homology};
\end{tikzpicture}

%------------ Homotopy invariance of homology ---------------
\begin{tikzpicture}
	\node [mybox] (box){%
		\begin{minipage}{0.3\textwidth}
			
		\end{minipage}
	};
	%------------ Homotopy invariance of homology Header ---------------------
	\node[fancytitle, right=10pt] at (box.north west) {Homotopy invariance of homology};
\end{tikzpicture}

%------------ Long exact sequence of homology of pairs ---------------
\begin{tikzpicture}
	\node [mybox] (box){%
		\begin{minipage}{0.3\textwidth}
			
		\end{minipage}
	};
	%------------ Long exact sequence of homology of pairs Header ---------------------
	\node[fancytitle, right=10pt] at (box.north west) {Long exact sequence of homology};
\end{tikzpicture}

%------------ Universal Coefficient theorem ---------------
\begin{tikzpicture}
	\node [mybox] (box){%
		\begin{minipage}{0.3\textwidth}
			
		\end{minipage}
	};
	%------------ Universal Coefficient theorem Header ---------------------
	\node[fancytitle, right=10pt] at (box.north west) {Universal Coefficient theorem};
\end{tikzpicture}

%------------ Kunneth theorem ---------------
\begin{tikzpicture}
	\node [mybox] (box){%
		\begin{minipage}{0.3\textwidth}
			
		\end{minipage}
	};
	%------------ Kunneth theorem Header ---------------------
	\node[fancytitle, right=10pt] at (box.north west) {Kunneth theorem};
\end{tikzpicture}

%------------ Singular cohomology ---------------
\begin{tikzpicture}
	\node [mybox] (box){%
		\begin{minipage}{0.3\textwidth}
			
		\end{minipage}
	};
	%------------ Singular cohomology Header ---------------------
	\node[fancytitle, right=10pt] at (box.north west) {Singular cohomology};
\end{tikzpicture}

%------------ Cohomology of pairs ---------------
\begin{tikzpicture}
	\node [mybox] (box){%
		\begin{minipage}{0.3\textwidth}
			
		\end{minipage}
	};
	%------------ Cohomology of pairs Header ---------------------
	\node[fancytitle, right=10pt] at (box.north west) {Cohomology of pairs};
\end{tikzpicture}

%------------ Cup product ---------------
\begin{tikzpicture}
	\node [mybox] (box){%
		\begin{minipage}{0.3\textwidth}
			
		\end{minipage}
	};
	%------------ Cup product Header ---------------------
	\node[fancytitle, right=10pt] at (box.north west) {Cup product};
\end{tikzpicture}

%------------ Poincare duality ---------------
\begin{tikzpicture}
	\node [mybox] (box){%
		\begin{minipage}{0.3\textwidth}
			
		\end{minipage}
	};
	%------------ Poincare duality Header ---------------------
	\node[fancytitle, right=10pt] at (box.north west) {Poincare duality};
\end{tikzpicture}
% ==============================================================
% Explicit examples
% ==============================================================

\end{multicols*}
\newpage
\begin{center}{\fontsize{35}{60}{\textcolor{DarkOrchid}{\textbf{Algebraic Topology Cheat Sheet: Explicit Examples}}}}\\
\end{center}
\begin{multicols*}{3}

%------------ Spheres and disks ---------------
\begin{tikzpicture}
	\node [mybox] (box){%
		\begin{minipage}{0.3\textwidth}
			
		\end{minipage}
	};
	%------------ Spheres and disks Header ---------------------
	\node[fancytitle, right=10pt] at (box.north west) {Spheres and disks};
\end{tikzpicture}

%------------ Graphs and wedges ---------------
\begin{tikzpicture}
	\node [mybox] (box){%
		\begin{minipage}{0.3\textwidth}
			
		\end{minipage}
	};
	%------------ Graphs and wedges Header ---------------------
	\node[fancytitle, right=10pt] at (box.north west) {Graphs and wedges};
\end{tikzpicture}

%------------ Products and smash ---------------
\begin{tikzpicture}
	\node [mybox] (box){%
		\begin{minipage}{0.3\textwidth}
			
		\end{minipage}
	};
	%------------ Products and smash Header ---------------------
	\node[fancytitle, right=10pt] at (box.north west) {Products and smash};
\end{tikzpicture}

%------------ Suspensions and cones ---------------
\begin{tikzpicture}
	\node [mybox] (box){%
		\begin{minipage}{0.3\textwidth}
			
		\end{minipage}
	};
	%------------ Suspensions and cones Header ---------------------
	\node[fancytitle, right=10pt] at (box.north west) {Suspensions and cones};
\end{tikzpicture}

%------------ Surfaces ---------------
\begin{tikzpicture}
	\node [mybox] (box){%
		\begin{minipage}{0.3\textwidth}
			
		\end{minipage}
	};
	%------------ Surfaces Header ---------------------
	\node[fancytitle, right=10pt] at (box.north west) {Surfaces};
\end{tikzpicture}

%------------ Projective spaces ---------------
\begin{tikzpicture}
	\node [mybox] (box){%
		\begin{minipage}{0.3\textwidth}
			
		\end{minipage}
	};
	%------------ Projective spaces Header ---------------------
	\node[fancytitle, right=10pt] at (box.north west) {Projective spaces};
\end{tikzpicture}

%------------ Lens spaces ---------------
\begin{tikzpicture}
	\node [mybox] (box){%
		\begin{minipage}{0.3\textwidth}
			
		\end{minipage}
	};
	%------------ Lens spaces Header ---------------------
	\node[fancytitle, right=10pt] at (box.north west) {Lens spaces};
\end{tikzpicture}

%------------ Grassmannians ---------------
\begin{tikzpicture}
	\node [mybox] (box){%
		\begin{minipage}{0.3\textwidth}
			
		\end{minipage}
	};
	%------------ Grassmannians Header ---------------------
	\node[fancytitle, right=10pt] at (box.north west) {Grassmannians};
\end{tikzpicture}

%------------ Flag manifolds ---------------
\begin{tikzpicture}
	\node [mybox] (box){%
		\begin{minipage}{0.3\textwidth}
			
		\end{minipage}
	};
	%------------ Flag manifolds Header ---------------------
	\node[fancytitle, right=10pt] at (box.north west) {Flag manifolds};
\end{tikzpicture}

%------------ Stiefel manifolds ---------------
\begin{tikzpicture}
	\node [mybox] (box){%
		\begin{minipage}{0.3\textwidth}
			
		\end{minipage}
	};
	%------------ Stiefel manifolds Header ---------------------
	\node[fancytitle, right=10pt] at (box.north west) {Stiefel manifolds};
\end{tikzpicture}

%------------ Classical groups ---------------
\begin{tikzpicture}
	\node [mybox] (box){%
		\begin{minipage}{0.3\textwidth}
			
		\end{minipage}
	};
	%------------ Classical groups Header ---------------------
	\node[fancytitle, right=10pt] at (box.north west) {Classical groups};
\end{tikzpicture}

%------------ Cayley projective plane ---------------
\begin{tikzpicture}
	\node [mybox] (box){%
		\begin{minipage}{0.3\textwidth}
			
		\end{minipage}
	};
	%------------ Cayley projective plane Header ---------------------
	\node[fancytitle, right=10pt] at (box.north west) {Cayley projective plane};
\end{tikzpicture}

%------------ Standard fibrations ---------------
\begin{tikzpicture}
	\node [mybox] (box){%
		\begin{minipage}{0.3\textwidth}
			
		\end{minipage}
	};
	%------------ Standard fibrations Header ---------------------
	\node[fancytitle, right=10pt] at (box.north west) {Standard fibrations};
\end{tikzpicture}

%------------ Standard homotopy groups ---------------
\begin{tikzpicture}
	\node [mybox] (box){%
		\begin{minipage}{0.3\textwidth}
			
		\end{minipage}
	};
	%------------ Standard homotopy groups Header ---------------------
	\node[fancytitle, right=10pt] at (box.north west) {Standard homotopy groups};
\end{tikzpicture}

%------------ Standard homology and cohomology ---------------
\begin{tikzpicture}
	\node [mybox] (box){%
		\begin{minipage}{0.3\textwidth}
			
		\end{minipage}
	};
	%------------ Standard homology and cohomology Header ---------------------
	\node[fancytitle, right=10pt] at (box.north west) {Standard homology and cohomology};
\end{tikzpicture}

\end{multicols*}
\end{document}